\section{Emission model: derivation, discrepancy, and reproduction}
\label{sec:app-emission}

Every emission figure in this paper is produced by one model, and this appendix states what that
model is, where its constants come from, one place where it disagrees with its own documentation,
and how to reproduce every number.

\subsection{Provenance}

The model is \texttt{supply/\allowbreak{}emission\_model.\allowbreak{}py} in the \texttt{flux-roadmap} repository. It is not a
reimplementation written for this paper: it is the team's own model, and it carries source citations
to \texttt{chainparams.\allowbreak{}cpp}, \texttt{main.cpp} and \texttt{amount.h} for each constant it uses. Those
citations were checked independently against the daemon source during the preparation of
\S\ref{sec:monetary-policy}, and \S\ref{sec:app-constants} reproduces the constants with their
locations.

\subsection{Arithmetic conventions, which are load-bearing}

All arithmetic is over integers in satoshi, matching \fluxd's \texttt{CAmount} (\texttt{int64}).
Division is floor division, matching C++ integer division for non-negative operands, and \texttt{>>}
is an arithmetic right shift. This is not pedantry: the \PoN reduction applies
$s_{k+1} = \lfloor 9 s_k / 10 \rfloor$ once per interval, and because each step truncates
independently, the result is \emph{not} in general equal to $\lfloor s_0 \cdot 0.9^{r} \rfloor$. Any
reimplementation using floating-point arithmetic will diverge from consensus.

\subsection{The two modes, and the discrepancy}

The model runs in two modes. \texttt{code} transcribes \texttt{Get\allowbreak{}Block\allowbreak{}Subsidy} literally.
\texttt{observed} reproduces what mainnet actually paid over blocks 1--4999, where the deployed
behaviour differs from a literal reading of the slow-start branch by one block of offset. The two
per-block subsidies converge for every height at or above 5{,}000.

The model's own docstring states that the two modes differ by \SI{0.06}{\flux}. Executing it shows an
exact and persistent difference in cumulative supply of \textbf{\num{150.00000000}\flux} at every height at or above
\num{5000}, arising from the one-block index shift itself, which lowers every ramp block's payout
rather than block 2's alone as the docstring's figure assumes. This is a
discrepancy in the published model's documentation, not in consensus: neither mode changes what the
chain paid, and the difference is confined to the first 5{,}000 blocks of a chain that is now past
2.9 million. This paper uses \texttt{observed} throughout, and reports supply as
\num{429466823.9700}\flux at height \num{2912015}.

\begin{remark}[Status of this discrepancy]
\label{rem:emission-discrepancy}
The \num{150.00000000}\flux{} figure and its root cause --- the index shift acting on every ramp
block, not only on block 2 --- were established by executing the published model, and are
reproducible by anyone who runs it \shipped
\srccode{flux-roadmap/\allowbreak{}supply/\allowbreak{}emission\_model.py}{207--229}
\srcmeasured{flux-roadmap/\allowbreak{}supply/\allowbreak{}emission\_model.py, both modes}{2026-09-03}. \textbf{This is
recorded, not outstanding, so far as this paper is concerned}: nothing here depends on the model's
docstring being right, every figure this paper reports comes from \texttt{observed} mode, and the
difference is confined to the first \num{5000} blocks of a chain now past \num{2.9} million. It has
no consensus consequence.

What remains is an upstream housekeeping item for the model's own maintainers rather than a gap in
this document: the docstring says \SI{0.06}{\flux} where execution gives \num{150}\flux, and the fix
is to the docstring's figure, since the closed-form sums agree with the per-block function in both
modes. It is recorded here so that a reader who runs the model and
sees a different number than the docstring promises knows why, and so that the correction does not
depend on this paper being read.
\end{remark}

\subsection{The functions}

\S\ref{sec:monetary-policy} states $\subsidypow$, $\subsidypon$ and $\supply$ in full. This appendix
adds only the two facts a reimplementer needs and the body omits: the model computes cumulative
supply as a sum of six independently-triggered step components rather than by iterating over blocks,
so it is exact at any height without a loop; and it treats the dev-fund remainder as a
\emph{remainder} of the reduced subsidy, never as a fixed \SI{0.5}{\flux}, so the dev fund declines
on the same schedule as the tier bases.

\subsection{Reproduction}

\begin{lstlisting}[caption={Reproducing every emission figure in this paper.},captionpos=b]
# Model report (mode=observed is the default)
python3 repos/flux-roadmap/supply/emission_model.py

# Regenerate the plotted data: emission.dat, tier_income.dat, pon_reductions.dat
python3 repos/flux-roadmap/supply/_gen_dat.py

# Live measurements (fleet, applications, pricing, value pools)
bash   scripts/04-measure.sh
python3 scripts/05-network-analysis.py

# PNR settlement analysis: drip trade-off and crossover series
python3 scripts/07-pnr-settlement.py
\end{lstlisting}

Every generated \texttt{.dat} file under \texttt{paper/data/} begins with \texttt{\#} comment lines
recording the generating command, the date, the source model, and the assumptions in force --- in
particular that future-height projections hold block spacing at the \SI{30}{\second} target, which
real spacing may not do.

\subsection{What the model does not do}

It does not model demand, revenue, or price. The \PNR crossover series in \S\ref{sec:pnr} and
\S\ref{sec:evaluation} are produced by a separate script under explicitly stated demand-growth
assumptions, with price exogenous and never forecast. The model also assumes the schedule as
currently written in consensus: it does not anticipate the \PNR activation, the parallel-asset cut,
or any future change to the reduction grid.

%% Drain this section's float queue before the next begins.
\FloatBarrier
